\begin{textsample}{POS Dim 3 – mg – Score 21.00 – mg/mg\_inf\_camp\_exp\_4.txt}  \label{ex:f3_pos_263}
2.6.2 Campos de Experiências : Corpo, \textbf{gestos} e movimentos Com o corpo ( por meio dos sentidos, \textbf{gestos}, movimentos impulsivos ou intencionais, coordenados intencionalmente ou de forma espontânea ), as \textbf{crianças} desde cedo exploram o mundo, o espaço e os objetos do seu entorno, estabelecem relações, expressam -se, \textbf{brincam} e produzem conhecimentos sobre si, sobre o outro, sobre o universo social e cultural, tornando -se, progressivamente, conscientes dessa corporeidade.
Por meio das diferentes linguagens, como a música, a dança, o teatro, as \textbf{brincadeiras} de faz de conta, elas se comunicam e se expressam no entrelaçamento entre corpo, emoção e linguagem.
As \textbf{crianças} conhecem e reconhecem as sensações e funções de seu corpo e, com seus \textbf{gestos} e movimentos, identificam suas potencialidades e seus limites, desenvolvendo, ao mesmo tempo, a consciência sobre o que é seguro e o que pode ser um risco à sua integridade física.
Na Educação \textbf{Infantil}, o corpo das \textbf{crianças} ganha \textbf{centralidade}, pois ele é o \textbf{partícipe} privilegiado das práticas pedagógicas de \textbf{cuidado} físico, orientadas para a emancipação e a liberdade, e não para a submissão.
Assim, a \textbf{instituição} escolar precisa promover oportunidades ricas para que as \textbf{crianças} possam, sempre animadas pelo espírito \textbf{lúdico} e na interação com seus \textbf{pares}, explorar e vivenciar um amplo repertório de movimentos, \textbf{gestos}, olhares, \textbf{sons} e \textbf{mímicas} com o corpo, para descobrir variados modos de ocupação e uso do espaço com o corpo ( tais como sentar com apoio, rastejar, engatinhar, escorregar, caminhar \textbf{apoiando} -se em berços, mesas e cordas, saltar, escalar, equilibrar -se, correr, dar cambalhotas, alongar -se etc. ). ( BRASIL, 2017, p.38-39 ).
O corpo, como estrutura biológica humana, precisa ser percebido na sua totalidade, conectado às práticas sociais, à cultura e ao conhecimento.
A corporeidade supera a \textbf{percepção} do corpo meramente biológico, se constitui na história de cada sujeito e extrapola a concepção do movimento apenas como deslocamento.
A corporeidade influencia os modos de ser, estar, sentir e se comunicar no mundo.
Para além da individualidade, constitui -se também na interação entre sujeitos em relação ao espaço.
O movimento humano consiste em uma forma de linguagem corporal em que a expressão de sentimentos, emoções e pensamentos está relacionada à cultura a qual pertence.
As \textbf{crianças}, desde o \textbf{nascimento}, movimentam -se e seus movimentos se aprimoram a cada \textbf{dia} por meio de experiências.
Desse modo, o trabalho com experiências corporais na Educação \textbf{Infantil} deve propiciar à \textbf{criança} a ampliação das possibilidades expressivas de seu corpo e o desenvolvimento de aspectos ligados à sua motricidade.
A música, a dança, o teatro e as \textbf{brincadeiras} de “ faz-de-conta ” são recursos pelos quais as \textbf{crianças} se comunicam e expressam suas emoções.
Cabe à \textbf{instituição} escolar promover oportunidades para que elas conheçam e vivenciem um amplo repertório de movimentos, \textbf{gestos}, olhares, \textbf{sons} e \textbf{mímicas} e descubram variados modos de ocupação e uso do espaço com o corpo. 2.6.3 Campo de experiências : Traços, \textbf{sons}, cores e formas Conviver com as diferentes manifestações artísticas, culturais e científicas, locais e universais, no cotidiano da \textbf{instituição} escolar, possibilita às \textbf{crianças}, por meio de experiências diversificadas, vivenciar diversas formas de expressão e linguagens, como as artes visuais, a música, o teatro, a dança e o audiovisual, entre outras.
Com base nessas experiências, elas se expressam por várias linguagens, criando suas próprias produções artísticas ou culturais, exercitando autoria ( coletiva e individual ) com \textbf{sons}, traços, \textbf{gestos}, danças, \textbf{mímicas}, encenações, canções, desenhos, modelagens, manipulação de diversos materiais e de recursos tecnológicos.
Essas experiências contribuem para que, desde muito \textbf{pequenas}, as \textbf{crianças} desenvolvam senso estético e crítico, o conhecimento de si mesmas, dos outros e da realidade que as cerca.
Portanto, a Educação \textbf{Infantil} precisa promover a participação das \textbf{crianças} em tempos e espaços para a produção, manifestação e apreciação artística, de modo a favorecer o desenvolvimento da sensibilidade, da criatividade e da expressão pessoal das \textbf{crianças}, permitindo que se apropriem e reconfigurem, permanentemente, a cultura e potencializem suas singularidades, ao ampliar repertórios e interpretar suas experiências e vivências artísticas. ( BRASIL, 2017, p.39 ).
As \textbf{instituições} de Educação \textbf{Infantil} devem se manter \textbf{atentas} em ofertar um ambiente acolhedor, com uma diversidade de objetos, materiais e espaços que estimulem a \textbf{curiosidade} e o senso explorador de cada \textbf{criança}.
Na atualidade, onde \textbf{sons} e imagens fazem parte do cotidiano, o senso estético deve ser continuamente aprimorado, incentivado pela intencionalidade do docente que necessita garantir a todas as \textbf{crianças} oportunidades para viver o prazer de pesquisar, experimentar um cenário com cores, \textbf{sons}, traços e formas marcantes, que traduzem a visualidade e a sonoridade presentes nas diferentes expressões artísticas.
A linguagem artística da \textbf{criança} é influenciada pela cultura, seja por meio de materiais com que realiza suas produções, por imagens, \textbf{sons} ou elementos de produção artística : desenho, ilustração, gravura, \textbf{pintura}, bordado, escultura, construção, fotografia, cinema, televisão, computação, dentre outras.
O importante é possibilitá -la viver experiências com a música, a \textbf{pintura}, a escultura e outras formas artísticas como a dança, a literatura e o teatro.
Essas modalidades integram aspectos sensíveis, afetivos, intuitivos, estéticos e comunicativos, contribuindo para o desenvolvimento da sua sensibilidade estética, criatividade e expressão pessoal, afirmando sua singularidade e reconstruindo a cultura.
As \textbf{crianças} têm suas próprias impressões, ideias e interpretações sobre a produção de arte e o fazer artístico, elaboradas a partir de suas experiências ao longo da vida.
As experiências vivenciadas nesse campo permitem que o mundo simbólico seja conhecido e ressignificado, desenvolvendo sua capacidade de criar, apreciar e sensibilizar -se, promovendo a sua formação integral.

% matched lemmas: apoiar, atentar, brincadeira, brincar, centralidade, criança, cuidado, curiosidade, dia, gesto, infantil, instituição, lúdico, mímico, nascimento, par, partícipe, pequeno, percepção, pintura, som
\end{textsample}
